\documentclass[11pt]{article}

\usepackage[T1]{fontenc}
\usepackage{lmodern}
\usepackage{amsmath,amssymb,amsthm}
\usepackage[margin=1in]{geometry}
\usepackage{booktabs}
\usepackage{enumitem}
\usepackage{microtype}
\usepackage[hidelinks]{hyperref}

\hypersetup{
  pdftitle={Far-left stability of a chemotaxis front: the sharp threshold chi=4},
  pdfauthor={ShenWork formalization project}
}

\newcommand{\R}{\mathbb{R}}
\newcommand{\norm}[1]{\left\lVert #1\right\rVert}
\newcommand{\abs}[1]{\left\lvert #1\right\rvert}
\newcommand{\code}[1]{\nolinkurl{#1}}
\setlength{\emergencystretch}{3em}
\sloppy

\newtheorem{theorem}{Theorem}
\newtheorem{proposition}{Proposition}
\newtheorem{lemma}{Lemma}
\newtheorem{remark}{Remark}

\title{Far-Left Stability of a Chemotaxis Front\\
\large The sharp threshold \(\chi=4\), a localized entropy proof, and the
necessary basin hypothesis}
\author{\code{ShenWork} formalization project}
\date{27 July 2026}

\begin{document}
\maketitle

\begin{abstract}
For the normalized one-dimensional chemotaxis--logistic equation
\[
 u_t=u_{zz}+(c-\chi v_z)u_z
      +u\bigl(1-\chi v+(\chi-1)u\bigr),\qquad -v_{zz}+v=u,
\]
we determine the exact stability threshold of the plateau \((u,v)=(1,1)\)
behind a traveling front.  The linear dispersion relation is strictly
negative for every Fourier mode exactly when \(0\le\chi<4\); at \(\chi=4\)
the mode \(k=1\) is neutral, and for \(\chi>4\) an unstable band appears.
A nonlinear relative-entropy calculation has the same sharp coefficient
\(1-\chi^2/16\).  We localize that calculation on the whole line with a
slowly varying weight and explicit boundary errors, and obtain uniform
far-left convergence throughout \(0\le\chi<4\), conditional on a persistent
positive far-left band and the stated derivative compactness.

The basin condition cannot be removed for the natural broad class of
front-like solutions.  A genuine traveling wave of speed \(s\), observed in
a faster frame \(c>s\), is positive at the left end on every time slice but
converges to zero at every fixed point of that frame.  This convective-escape
example is formalized for an actual PDE orbit already at \(\chi=0\).
All theorem and counterexample statements cited below are checked in Lean
4.30.0 / Mathlib v4.30.0, without \code{sorry}; their axiom audits use only
\code{propext}, \code{Classical.choice}, and \code{Quot.sound}.
\end{abstract}

\tableofcontents

\section{Problem and normalization}

In the traveling-wave problem the original population--signal system is
placed in the frame \(z=x-ct\).  In the normalized case
\[
  m=\alpha=\gamma=1
\]
the equation can be written in either of two equivalent ways:
\begin{align}
 u_t
  &=u_{zz}+(c-\chi v_z)u_z
    +u\bigl(1-\chi v+(\chi-1)u\bigr),
    &-v_{zz}+v&=u,                                      \label{eq:expanded}\\
 u_t
  &=u_{zz}+cu_z-\chi(uv_z)_z+u(1-u),
    &-r_{zz}+r&=u-1,\qquad r=v-1.                       \label{eq:divergence}
\end{align}
Indeed \(v_{zz}=v-u\), so expanding the divergence in
\eqref{eq:divergence} gives \eqref{eq:expanded}.

The traveling profile connects
\[
  (u,v)=(1,1)\quad\text{at }z=-\infty
  \qquad\text{to}\qquad
  (u,v)=(0,0)\quad\text{at }z=+\infty.
\]
The question addressed here is not construction of the wave, but stability
of the left plateau in the moving frame.  The target conclusion is the
quantified statement
\begin{equation}
 \forall\varepsilon>0\ \exists R,T\ \forall t\ge T\ \forall z\le-R,
 \qquad \abs{u(t,z+ct)-1}<\varepsilon.                  \label{eq:uniform-left}
\end{equation}
In Lean this is
\code{UniformCoMovingLeftEquilibriumConvergence c orbit}.

\section{The exact linear threshold}

\subsection{Dispersion relation}

Write
\[
  u=1+p,\qquad v=1+r.
\]
The elliptic equation linearizes exactly to
\[
  -r_{zz}+r=p.
\]
The product \(r_zp_z\) is quadratic, while the linearization of the reaction
term in \eqref{eq:expanded} is \((\chi-1)p-\chi r\).  Hence
\[
  p_t=p_{zz}+cp_z+(\chi-1)p-\chi r.
\]
For a Fourier mode \(p=e^{\lambda t+ikz}\), \(s=k^2\ge0\), the resolver gives
\[
  \widehat r=\frac{\widehat p}{1+s},
\]
and therefore
\begin{equation}
 \lambda(s)=-s+ick+(\chi-1)-\frac{\chi}{1+s}.
 \label{eq:lambda}
\end{equation}
The transport \(ick\) affects only the imaginary part.  The real growth rate
has the useful form
\begin{equation}
 d_\chi(s):=\operatorname{Re}\lambda(s)
   =-1-s+\frac{\chi s}{1+s}.                            \label{eq:dispersion}
\end{equation}

\subsection{Maximization}

There are two regimes.  If \(0\le\chi\le1\), then
\[
  d_\chi'(s)=-1+\frac{\chi}{(1+s)^2}\le0,
\]
so the maximum on \(s\ge0\) is \(d_\chi(0)=-1\).
If \(\chi>1\), the unique critical point is
\[
  s_*=\sqrt\chi-1>0,
\]
and it is the global maximizer.  Substitution gives
\begin{equation}
  d_\chi(s_*)=(\sqrt\chi-1)^2-1
             =\chi-2\sqrt\chi.                         \label{eq:max}
\end{equation}
Consequently
\[
  d_\chi(s)<0\quad\text{for all }s\ge0
  \quad\Longleftrightarrow\quad 0\le\chi<4.
\]
At \(\chi=4\), \(s_*=1\), hence \(k_*=\pm1\), and the maximal real part is
zero.  For \(\chi>4\), the maximizing mode has positive growth.

The formalized general scalar dispersion is
\[
  d_{\alpha,\chi\gamma}(s)
    =-\alpha-s+\frac{\chi\gamma\,s}{1+s}.
\]
For \(\chi\gamma>1\), its maximum is
\[
  (\sqrt{\chi\gamma}-1)^2-\alpha,
\]
so the exact threshold is
\begin{equation}
  \chi\gamma<(1+\sqrt\alpha)^2.                         \label{eq:general-threshold}
\end{equation}
For \(\chi\gamma\le1\), the maximum is the zero mode
\(-\alpha\), which is already negative.

\begin{theorem}[sharp dispersion threshold]
Let \(\alpha>0\) and \(\chi\gamma\ge0\).
If \(\chi\gamma<(1+\sqrt\alpha)^2\), then
\(d_{\alpha,\chi\gamma}(s)<0\) for every \(s\ge0\).
If \(\chi\gamma>(1+\sqrt\alpha)^2\), then some \(s\ge0\)
has positive growth.  When \(\chi\gamma>1\), the maximum is attained at
\(s=\sqrt{\chi\gamma}-1\).
\end{theorem}

\noindent The corresponding Lean theorems are
\code{dispersion_le_of_lt_turing},
\code{dispersion_pos_of_gt_turing}, and
\code{dispersion_attains_at_sqrt} in
\code{Paper1/WholeLineChiPosDispersionSharp.lean}.

\section{The sharp nonlinear entropy identity}

\subsection{Relative entropy and the exact cancellation}

For \(u>0\), define
\[
  h(u)=u-1-\log u,\qquad h'(u)=1-\frac1u.
\]
Let
\[
  A=\frac{u_z}{u},\qquad W=u-1.
\]
On a periodic interval, or before boundary terms are restored, multiply
\eqref{eq:divergence} by \(h'(u)\) and integrate.  The diffusion term is
\[
  \int h'(u)u_{zz}=-\int A^2.
\]
The constant drift is a total derivative and integrates to zero.  For the
chemotaxis term,
\begin{align*}
 -\chi\int h'(u)(uv_z)_z
 &= -\chi\int
     \left(1-\frac1u\right)(u_zr_z+ur_{zz})\\
 &= \chi\int A r_z.
\end{align*}
The second equality follows by integrating the \(Wr_{zz}\) term by parts;
this is the cancellation that keeps the logarithmic entropy sharp.  Finally,
\[
  \int h'(u)u(1-u)=-\int W^2.
\]
Thus
\begin{equation}
  \frac{d}{dt}\int h(u)
   =-\int A^2+\chi\int A r_z-\int W^2.                  \label{eq:entropy-raw}
\end{equation}

Young's inequality gives
\[
  \chi A r_z\le A^2+\frac{\chi^2}{4}r_z^2.
\]
The coefficient \(1/4\) in the one-dimensional elliptic resolver estimate is
also sharp:
\begin{equation}
  \int r_z^2\le\frac14\int W^2,
  \qquad -r_{zz}+r=W.                                  \label{eq:resolver}
\end{equation}
In Fourier variables, the multiplier is
\[
  \frac{k^2}{(1+k^2)^2}\le\frac14,
\]
with equality at \(k^2=1\).  Combining
\eqref{eq:entropy-raw}--\eqref{eq:resolver} yields
\begin{equation}
  \frac{d}{dt}\int h(u)
  \le-\left(1-\frac{\chi^2}{16}\right)\int(u-1)^2.
  \label{eq:sharp-entropy}
\end{equation}
The coefficient is positive exactly for \(\abs\chi<4\).  In the attractive
range considered here, this is precisely \(0\le\chi<4\), matching the linear
threshold and the critical mode.

\noindent Equation \eqref{eq:sharp-entropy} is formalized as
\code{sharp_entropy_production_le} in
\code{Paper1/WholeLineChiPosSharpEntropyDissipation.lean}.  The proof is
genuinely nonlinear: it assumes positivity of \(u\), not smallness of
\(u-1\).

\subsection{Why Fourier analysis is not enough on the moving far left}

The target \eqref{eq:uniform-left} is a whole-line, moving-half-line
statement.  A periodic calculation has no boundary flux, while a sharp
Fourier multiplier does not directly localize to a moving window.  The
formalization therefore proves the resolver estimate in real space with a
positive weight \(w\) satisfying
\begin{equation}
  \abs{w'(z)}\le\frac1\ell w(z),\qquad \ell>0.           \label{eq:slow-weight}
\end{equation}
On a finite interval \([a,b]\), integration by parts gives
\begin{align}
 \int_a^b w r_z^2
 &\le \frac14\int_a^b wW^2
 +\frac1{2\ell}
   \left(\int_a^b wr^2+\int_a^b wr_z^2\right) \notag\\
 &\quad
 +\abs{w(b)r(b)r_z(b)-w(a)r(a)r_z(a)}.                 \label{eq:weighted-resolver}
\end{align}
The leading source coefficient remains exactly \(1/4\); every loss is
displayed as a slow-weight or endpoint term.

\subsection{The localized production inequality}

Applying the same weighted integration by parts to the population equation
gives
\begin{align}
 \int_a^b w h'(u)u_t
 &\le
 -\left(1-\frac{\chi^2}{16}\right)
     \int_a^b wW^2                                      \label{eq:weighted-main}\\
 &\quad+\frac{\abs c+\chi+2}{\ell}
   \left[
     \int_a^b wh(u)+
     \int_a^b w(A^2+W^2)
   \right] \notag\\
 &\quad+
 \left(\frac{\chi^2}{4}+\frac{\chi}{2\ell}\right)
 \left[
  \frac1{2\ell}\left(\int_a^b wr^2+\int_a^b wr_z^2\right)
  +B_{\mathrm{res}}
 \right]
 +B_{\mathrm{pop}}, \notag
\end{align}
where the boundary terms are explicit:
\begin{align*}
 B_{\mathrm{res}}
  &=\abs{w(b)r(b)r_z(b)-w(a)r(a)r_z(a)},\\
 B_{\mathrm{pop}}
  &=\bigl[w h'(u)u_z\bigr]_a^b
    +c\bigl[wh(u)\bigr]_a^b
    -\chi\bigl[wWr_z\bigr]_a^b.
\end{align*}
No unspecified \(O(1/\ell)\) term hides dependence on the solution.  The
complete inequality is
\code{weighted_sharp_entropy_production_le}, and
\eqref{eq:weighted-resolver} is \code{weighted_resolver_le}, both in
\code{Paper1/WholeLineChiPosWeightedEntropyDissipation.lean}.  The
instantiation on the canonical Cauchy orbit is
\code{wholeLineCauchyGlobalU_weighted_sharp_dissipation}; the literal time
derivative of the entropy window is proved in
\code{Paper1/WholeLineChiPosWeightedEntropyOrbitDeriv.lean}.

\section{From entropy to far-left convergence}

\subsection{The two nontrivial hypotheses}

The theorem uses a persistent far-left band
\begin{equation}
 \exists T,Z\ \forall t\ge T\ \forall z\le Z,\qquad
   \ell\le u(t,z+ct)\le M,\qquad \ell>0,                \label{eq:band}
\end{equation}
formalized as \code{EventualCoMovingLeftBand c ell M orbit}.
The upper bound is available from the global ceiling.  The important part is
the time-uniform positive floor.

It also assumes \code{ChiOneFarLeftDerivativeCompactness}: every sequence of
time shifts \(t_n\to+\infty\) and far-left space shifts \(z_n\to-\infty\)
carries the first- and second-order population/resolver fields with the
uniform estimates needed to extract a classical translated limit.  This is
the precise compactness input; the theorem does not hide it inside a
statement called ``regularity.''

\subsection{Liouville/compactness mechanism}

The proof follows four steps.
\begin{enumerate}[leftmargin=2em]
\item Suppose \eqref{eq:uniform-left} fails.  Then some
  \(\varepsilon_0>0\) and sequences \(t_n\to\infty\),
  \(z_n\to-\infty\) satisfy
  \(\abs{u(t_n,z_n+ct_n)-1}\ge\varepsilon_0\).
\item Translate the orbit by \((t_n,z_n)\).  The band
  \eqref{eq:band} gives uniform positive lower and upper bounds on every
  fixed compact box.  Derivative compactness extracts an entire classical
  limit.
\item Apply the localized entropy inequality on successively larger windows.
  The slow-weight and endpoint errors vanish in the exhaustion.  Since
  \(1-\chi^2/16>0\), the entire limit has zero displacement dissipation:
  \(\int (u-1)^2=0\) on every compact interval.
\item Hence the limit is \(u\equiv1\), contradicting the nontrivial value at
  the translated origin.
\end{enumerate}

\begin{theorem}[basin-conditional sharp-range convergence]
Let \(0\le\chi<4\), \(\ell>0\), and let an orbit satisfy
\code{EventualCoMovingLeftBand c ell M orbit} and
\code{ChiOneFarLeftDerivativeCompactness chi c orbit}.  Then
\code{UniformCoMovingLeftEquilibriumConvergence c orbit}; equivalently,
\eqref{eq:uniform-left} holds.
\end{theorem}

\noindent This is
\code{uniformCoMovingLeftEquilibriumConvergence_chiPos_upto_four_basinConditional}
in
\code{Paper1/WholeLineChiPosEntropyFarLeftBasinConditional.lean}.  The same
file contains the canonical-orbit version.  A separate theorem derives the
band from a persistent lower-barrier plateau and the global ceiling, leaving
only the derivative compactness input.

\section{Why the basin floor cannot be omitted}

\subsection{The quantifier distinction}

Positivity at the left endpoint on each time slice says
\begin{equation}
 \forall t\ \exists\ell_t>0\ \exists Z_t\ \forall z\le Z_t,
 \qquad u(t,z+ct)\ge\ell_t.                             \label{eq:slice}
\end{equation}
The persistent band needed above says
\begin{equation}
 \exists\ell>0\ \exists T,Z\ \forall t\ge T\ \forall z\le Z,
 \qquad u(t,z+ct)\ge\ell.                              \label{eq:uniform}
\end{equation}
The order of the first three quantifiers is reversed.  Classical positivity
on fixed slices does not imply \eqref{eq:uniform}.

\subsection{Convective escape}

Let \((U,V)\) be a genuine traveling wave of laboratory speed \(s\):
\[
  u_{\mathrm{lab}}(t,x)=U(x-st),\qquad
  v_{\mathrm{lab}}(t,x)=V(x-st),
\]
with \(U(-\infty)=1\) and \(U(+\infty)=0\).  Observe it in a frame of speed
\(c>s\).  The population profile is
\begin{equation}
  q(t,z)=u_{\mathrm{lab}}(t,z+ct)
        =U\bigl(z+(c-s)t\bigr).                         \label{eq:escape}
\end{equation}
For every fixed \(t\),
\[
  \lim_{z\to-\infty}q(t,z)=1,
\]
so \eqref{eq:slice} holds.  On the other hand, for every fixed \(z\),
\[
  z+(c-s)t\longrightarrow+\infty,
  \qquad q(t,z)\longrightarrow0.
\]
Therefore, for every proposed \(\ell>0,T,Z\), one can choose \(t\ge T\)
so that \(q(t,Z)<\ell\).  No persistent floor exists.

This is a speed mismatch, not a failure of pointwise positivity and not an
artifact of a fabricated trajectory.  The chain rule verifies that
\((q,W)\), with \(W(t,z)=V(z+(c-s)t)\), solves the exact moving-frame PDE.
The repository constructs a concrete example at
\[
  \chi=0,\qquad s=3,\qquad c=4.
\]

\begin{proposition}[genuine PDE obstruction]
There exists a global classical traveling-wave orbit at \(\chi=0\) such that
every fixed co-moving time slice is strictly positive at the left endpoint,
but no positive \code{EventualCoMovingLeftBand} exists in the faster frame.
\end{proposition}

\noindent The pure quantifier theorem is
\code{convectiveEscapeProfile_no_uniform_left_floor}; the wave-level
statement is
\code{IsTravelingWave.convectiveEscape_quantifier_obstruction}; and the
constructed \(\chi=0\) PDE counterexample is
\code{exists_genuine_convectiveEscape_counterexample_chi_zero} in
\code{Paper1/WholeLineConvectiveEscapePDE.lean}.

\subsection{What the counterexample does and does not refute}

It refutes an unconditional theorem quantified over broad front-like data
and an arbitrary observation speed.  It does not refute the paper's
wave-stability theorem with a weighted initial-closeness condition to a fixed
wave.  Such a condition pins the orbit and the frame to the same wave speed;
the convective-escape orbit is not close to the incorrectly framed wave in
the required weighted norm.

Thus the optimal sharp-range statement has one of two equivalent kinds of
extra information:
\begin{itemize}[leftmargin=2em]
\item a persistent plateau-band hypothesis such as \eqref{eq:band}; or
\item a wave-closeness/phase condition strong enough to rule out speed
  mismatch and imply that band.
\end{itemize}

\section{Audited alternative routes}

Several plausible approaches were developed far enough to identify their
exact obstruction.  They are recorded here to distinguish a failed estimate
from an untried idea.

\paragraph{Localized \(L^1\) mass.}
For a cutoff of width \(L\), \(\abs{\phi_L'}=O(L^{-1})\), but the transition
annulus has length \(O(L)\).  With only a uniform bound on the orbit, the
transport/chemotaxis leakage is \(O(1)\), not \(o(1)\).  Increasing the window
therefore does not prevent all mass from convecting through its boundary.
The escape profile \eqref{eq:escape} realizes exactly this mechanism.

\paragraph{Bare \(L^2\) deficit energy.}
The linearized quadratic form has the correct gap
\(2\sqrt\chi-\chi\) in the range \(1<\chi<4\), but the nonlinear
intermediate-amplitude terms are not sign-definite; for \(\chi>2\), the raw
quadratic deficit energy can increase.  The formalized files
\code{WholeLineChiPosL2SpectralCoercivity.lean},
\code{WholeLineChiPosL2DeficitEnergy.lean}, and
\code{WholeLineChiPosL2FarLeftDecay.lean} isolate the valid shallow-basin
part.  The logarithmic entropy succeeds because its multiplier cancels the
nonlinear mobility exactly.

\paragraph{Deep-hole refill.}
At a sufficiently deep interior minimum, locality of the elliptic resolver
makes \(v\) small and the logistic reaction positive, uniformly in \(\chi\).
This KPP-type refill mechanism is formalized in
\code{WholeLineChiPosDeepHoleRefill.lean}.  It controls symmetric interior
holes, but it cannot stop a monotone transition layer from crossing the
boundary of the moving half-line.

\paragraph{Front anchoring.}
If convergence to the selected traveling wave were already known, the
far-left floor would follow.  Using that convergence to prove the floor is
circular in the enlarged \(\chi<4\) range; the available wave-stability
assembly itself consumes far-left control or weighted closeness.

\paragraph{Parabolic Harnack.}
A Harnack inequality can propagate a positive local-mass seed into a pointwise
floor.  It does not create the seed.  The missing estimate is precisely a
uniform late-time local mass bound in a fixed moving half-line, and convective
escape shows that no such seed exists without phase/speed information.

\section{Lean theorem map and final conclusion}

\begin{center}
\small
\renewcommand{\arraystretch}{1.18}
\begin{tabular}{@{}p{0.36\textwidth}p{0.55\textwidth}@{}}
\toprule
Mathematical component & Principal source file / theorem\\
\midrule
Sharp dispersion threshold &
\code{WholeLineChiPosDispersionSharp.lean}:
\code{dispersion_le_of_lt_turing},
\code{dispersion_pos_of_gt_turing}\\
Sharp periodic entropy &
\code{WholeLineChiPosSharpEntropyDissipation.lean}:
\code{sharp_entropy_production_le}\\
Weighted resolver and entropy &
\code{WholeLineChiPosWeightedEntropyDissipation.lean}:
\code{weighted_resolver_le},
\code{weighted_sharp_entropy_production_le}\\
Canonical-orbit instantiation &
\code{WholeLineChiPosWeightedEntropyOrbit.lean}:
\code{wholeLineCauchyGlobalU_weighted_sharp_dissipation}\\
Basin-conditional convergence &
\code{WholeLineChiPosEntropyFarLeftBasinConditional.lean}:
\code{uniformCoMovingLeftEquilibriumConvergence_chiPos_upto_four_basinConditional}\\
Convective quantifier obstruction &
\code{WholeLineChiPosFloorQuantifierObstruction.lean}:
\code{convectiveEscapeProfile_no_uniform_left_floor}\\
Genuine PDE counterexample &
\code{WholeLineConvectiveEscapePDE.lean}:
\code{exists_genuine_convectiveEscape_counterexample_chi_zero}\\
\bottomrule
\end{tabular}
\end{center}

\medskip
\noindent The value \(4\) is sharp in three independent senses:
\begin{enumerate}[leftmargin=2em]
\item below \(4\), every linear Fourier mode decays; at \(4\), the mode
  \(k^2=1\) is neutral; above \(4\), a mode grows;
\item the nonlinear entropy has the exact coercive coefficient
  \(1-\chi^2/16\);
\item throughout \(0\le\chi<4\), the localized entropy and compactness
  argument proves uniform far-left convergence once a persistent plateau
  band is supplied.
\end{enumerate}
The basin/phase condition is structural: without it, a genuine front can
escape convectively from a faster observation frame.  The resolved theorem is
therefore the basin-conditional sharp-range result, together with the
convective-escape counterexample delimiting any unconditional extension.

\end{document}
